\section{Table and annotation reference}
\begin{tabularx}{\linewidth}{@{}>{\raggedright\arraybackslash}p{55mm}X@{}}
\toprule
\multicolumn{1}{c}{Interface / option} & \multicolumn{1}{c}{Purpose}\\\midrule
\key{\gradcolumn[...]{id}{title}} & Declare a unique column before rows.\\
\key{\gradrow{cells}} & Comma-separated cells in column order.\\
\key{\gradheader{spans}} & Add a full header layer before data.\\
\key{\gradspan{n}{title}} & Span n consecutive columns in that layer.\\
\key{\gradgroup{title}} & Full-width heading followed by data.\\
\key{\gradtablestyle{name}{keys}} & Reuse table options with style=name.\\
\key{\gradtableschema{name}{body}} & Reuse declarations with graduseschema.\\
\key{mark=none|best|both} & Raw-value ranking for number/bar/bullet.\\
\key{width}, \key{weight} & Total width and flexible-column allocation.\\
\key{long=true} & Paginate, repeating complete headers.\\
\key{caption}, \key{label} & Long-table title and reference; label needs caption.\\
\key{overflow mark} & none, arrow or slash at clipped boundaries.\\
\key{marker shape} & circle, square, diamond or triangle.\\
\key{compare annotation} & none, arrow or bracket.\\
\key{compare label} & values, delta or percent.\\
\key{references} & value/color/style/name records.\\
\key{bullet bands} & Increasing upper-bound/color records ending at max.\\
\bottomrule
\end{tabularx}
Neither table buffering nor the drawing commands infer scientific meaning from
your data. Set scales, units, uncertainty interpretation, category identity and
quality direction explicitly. Horizontal header spans are supported; vertical
cell merging is not supported. Scalar overall and group summaries are supported.

\subsection{Tagged PDF and alternative text}
On a LaTeX distribution supporting PDF tagging, enable it with DocumentMetadata
before documentclass. Independent graphics generate English alternative text from
raw data: values, units, scales, error distances, reference lines, comparison
endpoints and differences, stack segments and subtotals, and every spark observation
including missing positions. Intervals describe supplied endpoints and estimates;
confidence levels are not inferred. Legends describe their category names.

\begin{source}
\DocumentMetadata{tagging=on,pdfstandard=ua-2,lang=en-US}
\documentclass{article}
\usepackage{gradbars}
\begin{document}
\gradbar{72}
\gradbar[alt={Accuracy: 72 percent, out of 100 percent.}]{72}
\graddumbbell[compare label=delta]{58}{76}
\end{document}
\end{source}
Use alt to supply context or another language. Empty alt restores automatic text;
the existing /tikz/alt syntax remains supported. Accessibility defaults to auto.
Use accessibility=artifact only when surrounding text fully conveys the graphic's
information. Accessibility=off leaves tagging to the surrounding code. These options
do not change visible output when tagging is disabled.

Descriptions use raw values before rounding. The text option only changes the
visual label; use alt to change its semantic description. Table measurement does
not produce tags, so each rendered graphic cell has one description. Inside an
existing TikZ picture, describe the complete graphic using the outer picture's alt.
Automatic descriptions do not establish document-wide PDF/UA conformance: validate
table headers, reading order, fonts and other content separately.

\subsection*{Contact and maintenance}
Author and maintainer: SuFan (Tape5try).
Contact \href{mailto:3546236610@qq.com}{3546236610@qq.com} or the
\href{https://github.com/Tape5try/gradbars/issues}{issue tracker}.
Include a minimal source, engine and log. Distributed under the MIT license.
